% ===========================================================================
%  Raport Experimental — Fazele 1 și 2
%  Cuantizare Numerică FP16 și INT8 pe ViT-Tiny pretrained (ImageNet-1k)
%  Tudose Alexandru · Lucrare de Licență 2026
% ===========================================================================
\documentclass[12pt,a4paper]{article}

% --- Encoding & fonts (XeTeX / tectonic) ---
\usepackage{fontspec}
\defaultfontfeatures{Ligatures=TeX}

% --- Language ---
\usepackage{polyglossia}
\setmainlanguage{romanian}

% --- Page geometry ---
\usepackage[
  top=2.5cm, bottom=2.5cm,
  left=2.5cm, right=2.5cm
]{geometry}

% --- Math ---
\usepackage{amsmath, amssymb}

% --- Graphics ---
\usepackage{graphicx}
\usepackage{float}
\usepackage{caption}
\usepackage{subcaption}
\captionsetup{font=small, labelfont=bf}

% --- Tables ---
\usepackage{booktabs}
\usepackage{array}
\usepackage{multirow}

% --- Colors ---
\usepackage{xcolor}
\definecolor{fp32color}{HTML}{0072B2}
\definecolor{fp16color}{HTML}{009E73}
\definecolor{int8color}{HTML}{E69F00}
\definecolor{lightgray}{HTML}{F4F4F4}

% --- Hyperlinks ---
\usepackage[
  colorlinks=true,
  linkcolor=fp32color,
  citecolor=fp32color,
  urlcolor=fp32color
]{hyperref}

% --- Misc ---
\usepackage{microtype}
\usepackage{parskip}
\setlength{\parindent}{0pt}
\setlength{\parskip}{6pt}
\usepackage{enumitem}

% --- Header / Footer ---
\usepackage{fancyhdr}
\pagestyle{fancy}
\fancyhf{}
\lhead{\small Tudose Alexandru --- Raport Experimental Faze 1--2}
\rhead{\small 2026}
\cfoot{\thepage}
\renewcommand{\headrulewidth}{0.4pt}

% ===========================================================================
\begin{document}
% ===========================================================================

% --- Title page ---
\begin{titlepage}
  \centering
  {\large Universitatea Transilvania din Brașov\\}
  {\small Facultatea de Matematică și Informatică\\}
  {\small Specializarea: Informatică Aplicată\\[2cm]}

  \rule{\linewidth}{0.4pt}\\[0.5cm]
  {\LARGE\bfseries Raport Experimental\\[0.3cm]}
  {\Large\bfseries Cuantizare Numerică FP16 și INT8\\[0.2cm]}
  {\Large\bfseries pe Vision Transformer Pretrained\\[0.2cm]}
  {\large (ImageNet-1k / ImageNette)}\\[0.3cm]
  \rule{\linewidth}{0.4pt}\\[2cm]

  \begin{minipage}{0.45\textwidth}
    \textbf{Student:}\\
    Tudose Alexandru\\
    Grupa 10LF333\\
    An III, Informatică Aplicată
  \end{minipage}
  \hfill
  \begin{minipage}{0.45\textwidth}
    \raggedleft
    \textbf{Coordonator:}\\
    Conf.\ dr.\ ing.\ Honorius Gâlmeanu\\
    Departamentul de Informatică
  \end{minipage}\\[3cm]

  {\small Brașov, 2026}
\end{titlepage}

% --- TOC ---
\tableofcontents
\newpage

% ===========================================================================
\section{Introducere și Motivație}
% ===========================================================================

Prezentul raport documentează rezultatele experimentelor din \textbf{Fazele 1 și 2}
ale lucrării de licență, ce studiază impactul cuantizării numerice asupra unui
Vision Transformer (ViT). Spre deosebire de studiul preliminar (antrenare pe CIFAR-10),
aceste faze utilizează un model \textbf{pretrained pe ImageNet-1k} și evaluează
cuantizarea exclusiv la inferență (\emph{post-training quantization, PTQ}), fără
nicio re-antrenare.

\textbf{Corecții față de studiul preliminar:}
\begin{itemize}
  \item Implementarea anterioară de FP8 folosea cast direct \texttt{torch.float8\_e4m3fn}
    — problematică deoarece nu toate valorile FP32 sunt reprezentabile. Faza 2
    implementează cuantizarea corectă: scalare liniară cu \texttt{int8}, garantând
    o mapare bijectivă la 256 valori discrete.
  \item LayerNorm și classification head sunt excluse de la cuantizare (sensibile la
    erori numerice).
  \item Modelul este acum evaluat pe \textbf{ImageNette2-320}, un subset de 10 clase
    din ImageNet-1k, cu ponderile originale pretrained — fără fine-tuning.
\end{itemize}

\textbf{Întrebări de cercetare:}
\begin{enumerate}
  \item Cuantizarea statică FP16 (\texttt{model.half()}) degradează acuratețea?
  \item Cuantizarea INT8 weight-only per-tensor introduce degradare semnificativă?
  \item Ce layere sunt cele mai sensibile la cuantizare (MSE per layer)?
  \item Care sunt câștigurile reale de memorie și latență pe Apple M4 (MPS)?
\end{enumerate}

% ===========================================================================
\section{Arhitectura Modelului}
% ===========================================================================

\subsection{Vision Transformer — ViT-Tiny}

Modelul utilizat este \texttt{vit\_tiny\_patch16\_224} din biblioteca \texttt{timm},
pretrained pe ImageNet-1k cu 5.7 milioane de parametri. Arhitectura include:

\begin{table}[H]
  \centering
  \caption{Parametrii arhitecturali ai ViT-Tiny}
  \begin{tabular}{ll}
    \toprule
    \textbf{Proprietate} & \textbf{Valoare} \\
    \midrule
    Parametri totali    & 5.7 milioane \\
    Patch size          & $16 \times 16$ pixeli \\
    Input resolution    & $224 \times 224$ \\
    Embedding dim       & 192 \\
    Număr blocuri       & 12 \\
    Număr capete (MHA)  & 3 \\
    Pre-training dataset & ImageNet-1k (1.2M imagini, 1000 clase) \\
    \bottomrule
  \end{tabular}
  \label{tab:vit_config}
\end{table}

Blocul central Transformer operează cu Multi-Head Self-Attention:
\begin{equation}
  \mathrm{Attention}(Q, K, V) = \mathrm{softmax}\!\left(\frac{QK^T}{\sqrt{d_k}}\right) V
\end{equation}
unde $Q, K, V \in \mathbb{R}^{N \times d_k}$ sunt proiecțiile query, key și value.
Fiecare bloc include și o rețea feed-forward (MLP):
\begin{equation}
  \mathrm{FFN}(x) = \mathrm{GELU}(xW_1 + b_1)W_2 + b_2
\end{equation}

Ambele sub-straturi sunt înconjurate de conexiuni reziduale și \textbf{LayerNorm}
--- acesta din urmă \emph{nu} este cuantizat în Faza 2 (sensibil la erori numerice).

\subsection{Straturile cuantizabile}

Din totalul de 48 straturi \texttt{nn.Linear} identificate:
\begin{itemize}
  \item \texttt{blocks.N.attn.qkv} (12 straturi) --- proiecție Q, K, V combinată
  \item \texttt{blocks.N.attn.proj} (12 straturi) --- proiecție output attention
  \item \texttt{blocks.N.mlp.fc1} (12 straturi) --- prima proiecție MLP
  \item \texttt{blocks.N.mlp.fc2} (12 straturi) --- a doua proiecție MLP
\end{itemize}

Excluse: \texttt{norm} (LayerNorm), \texttt{cls\_token}, \texttt{pos\_embed},
\texttt{head} (classification head).

\subsection{Preprocessing corect}

\textbf{Notă importantă:} \texttt{timm} folosește pentru \texttt{vit\_tiny\_patch16\_224}
normalizare cu $\mu = \sigma = 0.5$ (nu normalizarea ImageNet standard
$\mu=[0.485, 0.456, 0.406]$). Utilizarea normalizării greșite a redus acuratețea de la
$\approx 79\%$ la $47\%$. Configul corect a fost obținut via:
\begin{verbatim}
data_config = timm.data.resolve_data_config({}, model=model)
val_transform = timm.data.create_transform(**data_config, is_training=False)
\end{verbatim}

% ===========================================================================
\section{Dataset și Configurație Experimentală}
% ===========================================================================

\subsection{ImageNette2-320}

ImageNette~\cite{imagenette} este un subset de 10 clase din ImageNet-1k, ales ca
proxy pentru evaluare rapidă locală (fără a necesita descărcarea celor 150 GB ai
ImageNet complet).

\begin{table}[H]
  \centering
  \caption{Configurație experiment}
  \begin{tabular}{ll}
    \toprule
    \textbf{Parametru} & \textbf{Valoare} \\
    \midrule
    Model               & \texttt{vit\_tiny\_patch16\_224} \\
    Pre-training        & ImageNet-1k (\texttt{pretrained=True}) \\
    Dataset evaluare    & ImageNette2-320 (10 clase, 3925 imagini val) \\
    Device              & Apple M4 MPS \\
    Batch size          & 64 \\
    PyTorch             & 2.9.1 \\
    \bottomrule
  \end{tabular}
  \label{tab:exp_config}
\end{table}

\textbf{Maparea label-urilor:} Modelul pretrained produce logits pentru 1000 clase
ImageNet-1k. Folderele ImageNette (WordNet IDs) sunt mapate la indicii corecți din
ImageNet-1k (ex. \texttt{n01440764} $\to$ clasa 0 ``tench'').

% ===========================================================================
\section{Faza 1 — Cuantizare FP16 Statică}
% ===========================================================================

\subsection{Metodologie}

Faza 1 evaluează conversia completă a ponderilor la \texttt{float16} prin
\texttt{model.half()}, urmată de inferență pe MPS:

\begin{verbatim}
model_fp16 = timm.create_model("vit_tiny_patch16_224",
                                pretrained=True, num_classes=1000)
model_fp16 = model_fp16.half().to(device).eval()
\end{verbatim}

Aceasta este o cuantizare \textbf{statică} (nu AMP/autocast) --- \emph{toate}
ponderile sunt stocate și procesate în \texttt{float16}, inclusiv LayerNorm.

\subsection{Rezultate}

\begin{table}[H]
  \centering
  \caption{Rezultate Faza 1 --- FP32 vs.\ FP16 pe ImageNette}
  \begin{tabular}{lrrrr}
    \toprule
    \textbf{Metodă} & \textbf{Acuratețe} & \textbf{Loss} &
    \textbf{Latență (ms/batch)} & \textbf{Memorie (MB)} \\
    \midrule
    FP32 (baseline) & 79.11\% & 0.7474 & 141.45 & 21.81 \\
    FP16 (\texttt{model.half()}) & 79.16\% & 0.7472 & 110.25 & 10.91 \\
    \midrule
    \textbf{Diferență} & $-0.051$ pp & --- & $\mathbf{1.28\times}$ mai rapid &
    $\mathbf{2.00\times}$ mai mic \\
    \bottomrule
  \end{tabular}
  \label{tab:fp16_results}
\end{table}

\begin{figure}[H]
  \centering
  \includegraphics[width=\textwidth]{../results/FP16ImageNet/plots/00_summary.png}
  \caption{Overview Faza 1: comparație FP32 vs.\ FP16 pe cele trei metrici principale.
    Acuratețea este practic identică ($\Delta = -0.051$ pp), memoria se reduce exact 2×
    (4 bytes $\to$ 2 bytes per parametru), latența scade cu 1.28×.}
  \label{fig:fp16_summary}
\end{figure}

\begin{figure}[H]
  \centering
  \begin{subfigure}[b]{0.48\textwidth}
    \includegraphics[width=\textwidth]{../results/FP16ImageNet/plots/01_accuracy_comparison.png}
    \caption{Acuratețe top-1}
  \end{subfigure}
  \hfill
  \begin{subfigure}[b]{0.48\textwidth}
    \includegraphics[width=\textwidth]{../results/FP16ImageNet/plots/03_latency_comparison.png}
    \caption{Latență inferență}
  \end{subfigure}
  \caption{Detaliu acuratețe și latență. FP16 este marginal \emph{mai bun} cu $-0.051$~pp
    (zgomot numeric, nu câștig real), și $1.28\times$ mai rapid pe MPS Apple M4.}
  \label{fig:fp16_detail}
\end{figure}

\subsection{Discuție}

\textbf{Acuratețe:} Degradarea de $-0.051$~pp (FP16 marginal mai bun) se încadrează
în variabilitatea numerică și nu poate fi atribuită preciziei. ViT-Tiny este robust
la conversia float32 $\to$ float16.

\textbf{Memorie:} Reducerea exactă de $2\times$ (21.81 MB $\to$ 10.91 MB) confirmă
teoria: \texttt{float16} ocupă exact jumătate față de \texttt{float32}.

\textbf{Latență:} Speedup-ul de $1.28\times$ pe Apple M4 MPS este real, dar mai modest
decât pe GPU NVIDIA (unde Tensor Cores pentru FP16 oferă $1.5\text{--}3\times$).
Apple Silicon valorifică \texttt{float16} nativ, dar fără unități dedicate.

\textbf{Compatibilitate MPS:} \texttt{model.half()} funcționează complet pe MPS cu
PyTorch 2.9, inclusiv LayerNorm în \texttt{float16} --- nu a fost necesar fallback pe CPU.

% ===========================================================================
\section{Faza 2 — Cuantizare INT8 per-tensor (weight-only)}
% ===========================================================================

\subsection{Metodologie}

Faza 2 implementează cuantizarea corectă pe 8 biți cu scalare liniară
(\emph{weight-only, per-tensor}):

\begin{equation}
  s = \frac{\max(|W|)}{127}, \quad
  q = \mathrm{clamp}\!\left(\mathrm{round}\!\left(\frac{W}{s}\right), -128, 127\right)
  \in \mathbb{Z}_{8}
\end{equation}

La inferență, ponderile sunt decuantizate înainte de multiplicare matriceală:
\begin{equation}
  \hat{W} = q \cdot s \in \mathbb{R}
\end{equation}

Implementarea utilizează un wrapper \texttt{QuantizedLinear} care stochează
\texttt{(weight\_int8, scale)} și dequantizează la fiecare forward pass.

\subsection{Straturi cuantizate}

\begin{table}[H]
  \centering
  \caption{Strategia de cuantizare selectivă}
  \begin{tabular}{lll}
    \toprule
    \textbf{Tip strat} & \textbf{Cuantizat?} & \textbf{Motivație} \\
    \midrule
    \texttt{blocks.N.attn.qkv}  & Da & matmul, cuantizabil \\
    \texttt{blocks.N.attn.proj} & Da & matmul, cuantizabil \\
    \texttt{blocks.N.mlp.fc1}   & Da & matmul, cuantizabil \\
    \texttt{blocks.N.mlp.fc2}   & Da & matmul, cuantizabil \\
    LayerNorm (\texttt{norm})   & \textbf{Nu} & sensibil la erori numerice \\
    \texttt{head} (classifier)  & \textbf{Nu} & strat final, impact direct \\
    \texttt{cls\_token}, \texttt{pos\_embed} & \textbf{Nu} & embeddings, nu Linear \\
    \bottomrule
  \end{tabular}
  \label{tab:quant_strategy}
\end{table}

\textbf{Total:} 48 straturi \texttt{nn.Linear} cuantizate, 1 strat exclus (\texttt{head}).

\subsection{Rezultate}

\begin{table}[H]
  \centering
  \caption{Rezultate Faza 2 --- FP32 vs.\ INT8 pe ImageNette}
  \begin{tabular}{lrrrr}
    \toprule
    \textbf{Metodă} & \textbf{Acuratețe} & \textbf{Loss} &
    \textbf{Latență (ms/batch)} & \textbf{Memorie (MB)} \\
    \midrule
    FP32 (baseline) & 79.11\% & 0.7474 & 140.16 & 21.81 \\
    INT8 (per-tensor) & 78.47\% & 0.7712 & 141.12 & \phantom{0}6.62 \\
    \midrule
    \textbf{Diferență} & $+0.637$ pp & $+0.024$ & $\approx 1.00\times$ &
    $\mathbf{3.29\times}$ mai mic \\
    \bottomrule
  \end{tabular}
  \label{tab:int8_results}
\end{table}

\begin{figure}[H]
  \centering
  \includegraphics[width=\textwidth]{../results/INT8ImageNet/plots/00_summary.png}
  \caption{Overview Faza 2: comparație FP32 vs.\ INT8. Degradarea de $+0.637$~pp
    este acceptabilă pentru weight-only INT8. Reducerea de memorie este $3.29\times$
    (biasul rămâne în \texttt{float32}). Nu există speedup real deoarece dequantizarea
    la fiecare forward pass anulează câștigul de bandwidth.}
  \label{fig:int8_summary}
\end{figure}

\subsection{Analiză per-layer}

\begin{figure}[H]
  \centering
  \includegraphics[width=\textwidth]{../results/INT8ImageNet/plots/01_per_layer_mse.png}
  \caption{MSE al erorii de cuantizare per strat (eroarea $\|W - \hat{W}\|_F^2 / n$).
    Cele mai multe straturi au MSE în intervalul $[3\times10^{-7}, 5\times10^{-6}]$.
    Straturile \texttt{blocks.7.mlp.fc1} și \texttt{blocks.7.mlp.fc2} sunt
    \textbf{outlieri} cu MSE de $\approx 40\times$ mai mare decât media --- cauzat
    de ponderi cu valori extreme (scale factor $\approx 0.02$, față de media
    $\approx 0.004$). Acestea sunt candidați pentru cuantizare per-channel în Faza 3.}
  \label{fig:int8_mse}
\end{figure}

\begin{figure}[H]
  \centering
  \includegraphics[width=\textwidth]{../results/INT8ImageNet/plots/02_per_layer_scale.png}
  \caption{\textit{Stânga}: Factorul de scalare $s = \max(|W|)/127$ pentru fiecare
    strat, în ordinea din model. Variabilitate ridicată în straturile finale ale blocului 7.
    \textit{Dreapta}: Distribuție boxplot per tip de strat. Straturile \texttt{mlp.fc1}
    și \texttt{mlp.fc2} prezintă cea mai mare variabilitate a scale-ului, indicând
    că granularitatea per-tensor este insuficientă pentru acestea.}
  \label{fig:int8_scale}
\end{figure}

\subsection{Motivul lipsei speedup-ului}

Deși ponderile sunt stocate în \texttt{int8} (1 byte), dequantizarea la
\texttt{float32} înainte de fiecare matmul introduce un overhead care
compensează câștigul de bandwidth. Speedup-ul real apare doar cu:
\begin{itemize}
  \item Kernel-uri dedicate INT8 GEMM (ex. \texttt{torch.int\_mm} pe CUDA)
  \item Hardware cu suport nativ INT8 (NVIDIA A100, H100 -- Tensor Cores)
  \item Activări cuantizate simultan (\emph{weight + activation quantization})
\end{itemize}
Pe MPS Apple M4, nu există suport nativ INT8 GEMM în PyTorch 2.9.

% ===========================================================================
\section{Comparație Globală: FP32 / FP16 / INT8}
% ===========================================================================

\begin{figure}[H]
  \centering
  \includegraphics[width=\textwidth]{../results/INT8ImageNet/plots/03_cross_phase_comparison.png}
  \caption{Comparație directă FP32 / FP16 / INT8 pe cele trei metrici: acuratețe,
    memorie și latență. FP16 domină: memorie $2\times$ mai mică și latență $1.28\times$
    mai mică, cu degradare nulă. INT8 oferă cea mai mare compresie de memorie ($3.29\times$),
    dar fără beneficiu de viteză pe acest hardware.}
  \label{fig:cross_comparison}
\end{figure}

\begin{figure}[H]
  \centering
  \includegraphics[width=\textwidth]{../results/INT8ImageNet/plots/04_tradeoff.png}
  \caption{\textit{Stânga}: Scatter accuracy vs.\ memorie. INT8 oferă cel mai bun raport
    memorie/acuratețe. \textit{Dreapta}: Scatter accuracy vs.\ latență. FP16 domină
    clar: acuratețe identică cu FP32, dar latență mai mică.}
  \label{fig:tradeoff}
\end{figure}

\begin{table}[H]
  \centering
  \caption{Tabel comparativ complet FP32 / FP16 / INT8}
  \begin{tabular}{lrrrrrr}
    \toprule
    \textbf{Metodă} & \textbf{Acc.} & \textbf{$\Delta$ Acc.} &
    \textbf{Loss} & \textbf{Latență} & \textbf{Mem.} & \textbf{Compresie} \\
    \midrule
    FP32        & 79.11\% & ---        & 0.7474 & 141 ms & 21.81 MB & $1\times$ \\
    FP16        & 79.16\% & $-0.05$ pp & 0.7472 & 110 ms & 10.91 MB & $2.00\times$ \\
    INT8 (w-o)  & 78.47\% & $+0.64$ pp & 0.7712 & 141 ms & \phantom{0}6.62 MB & $3.29\times$ \\
    \bottomrule
  \end{tabular}
  \label{tab:global_comparison}
\end{table}

\textit{w-o = weight-only; latența este per batch (bs=64) pe Apple M4 MPS.}

% ===========================================================================
\section{Analiză și Discuții}
% ===========================================================================

\subsection{FP16 --- conversie fără cost}

Conversia \texttt{float32} $\to$ \texttt{float16} nu introduce nicio degradare
practică de acuratețe ($-0.051$~pp, în limitele zgomotului numeric). Câștigurile sunt:
\begin{itemize}
  \item \textbf{Memorie:} $2\times$ reducere exactă, conformă cu teoria
    (4 bytes $\to$ 2 bytes per parametru)
  \item \textbf{Latență:} $1.28\times$ speedup real pe MPS (Apple M4 valorifică
    \texttt{float16} nativ, fără unități Tensor Core dedicate)
  \item \textbf{Compatibilitate MPS completă:} \texttt{model.half()} funcționează
    fără excepții pe PyTorch 2.9, inclusiv LayerNorm
\end{itemize}

\subsection{INT8 weight-only --- compresie maximă fără speedup}

Cuantizarea INT8 per-tensor oferă cea mai mare compresie de memorie ($3.29\times$),
cu o degradare de acuratețe acceptabilă ($+0.637$~pp). Însă \textbf{nu există speedup}
față de FP32 deoarece implementarea dequantizează la \texttt{float32} la fiecare
forward pass.

Outlierul identificat (\texttt{blocks.7.mlp.*}, MSE $\approx 40\times$ media)
sugerează că o strategie \emph{per-channel} ar reduce semnificativ eroarea de
cuantizare pe aceste straturi specifice --- aceasta este una din analizele planificate
pentru Faza 3.

\subsection{Implicații pentru deployment}

\begin{itemize}
  \item \textbf{Pe Apple M4 (MPS):} FP16 este metoda optimă --- fără pierdere de
    acuratețe, $1.28\times$ mai rapid, $2\times$ mai puțin memorie.
  \item \textbf{Pe NVIDIA (CUDA):} INT8 cu kernel-uri GEMM native (Tensor Cores)
    ar oferi atât compresie de memorie ($3\text{--}4\times$) cât și speedup real
    ($2\text{--}4\times$) față de FP32.
  \item \textbf{FP16 vs.\ INT8:} FP16 câștigă pe acuratețe și latență (pe hardware
    cu suport nativ), INT8 câștigă pe memorie --- decizia depinde de constrângerea
    dominantă a aplicației.
\end{itemize}

% ===========================================================================
\section{Limitări ale studiului}
% ===========================================================================

\begin{itemize}
  \item \textbf{Dataset:} ImageNette are doar 10 clase ``ușoare'' --- rezultatele
    pot să nu se generalizeze la sarcini mai dificile (detecție, segmentare).
  \item \textbf{Hardware:} Apple M4 MPS nu are suport INT8 GEMM nativ --- concluzia
    privind lipsa speedup-ului INT8 este specifică acestei arhitecturi.
  \item \textbf{Weight-only:} Activările rămân în \texttt{float32} --- cuantizarea
    completă (weights + activations) ar putea oferi compresia deplină așteptată teoretic.
  \item \textbf{Per-tensor:} Granularitatea per-tensor este suboptimală pentru
    straturile cu distribuții de ponderi non-uniforme (\texttt{blocks.7.mlp.*}).
\end{itemize}

% ===========================================================================
\section{Concluzii și Pași Următori}
% ===========================================================================

\subsection{Concluzii}

\begin{enumerate}
  \item \textbf{FP32 $\to$ FP16 (static):} Degradare nulă ($-0.051$~pp),
    $2\times$ mai puțin memorie, $1.28\times$ mai rapid pe MPS.
    \textbf{Recomandare: utilizat implicit pentru deployment pe Apple Silicon.}

  \item \textbf{FP32 $\to$ INT8 (weight-only, per-tensor):} Degradare acceptabilă
    ($+0.637$~pp), $3.29\times$ mai puțin memorie, fără speedup real pe MPS din
    cauza overhead-ului de dequantizare.

  \item \textbf{Outlieri identificați:} Straturile \texttt{blocks.7.mlp.fc1} și
    \texttt{blocks.7.mlp.fc2} au MSE de cuantizare $\approx 40\times$ mai mare
    decât media --- candidați prioritari pentru per-channel quantization.

  \item \textbf{Preprocessing \texttt{timm}:} Normalizarea corectă trebuie obținută
    din \texttt{timm.data.resolve\_data\_config} --- nu din constante hardcoded
    (diferența: 79\% vs.\ 47\% acuratețe cu normalizare greșită).
\end{enumerate}

\subsection{Faza 3 --- Experimente planificate}

\begin{itemize}
  \item \textbf{Per-channel vs.\ per-tensor INT8} --- compararea granularităților
  \item \textbf{Sensitivity analysis} --- cuantizarea câte unui layer pe rând
    (impact individual per strat)
  \item \textbf{Layer-wise error} --- MSE/MAE per strat, vizualizare heatmap
  \item \textbf{Memory footprint} --- dimensiunea modelului serializat în fiecare format
    (\texttt{.pt}, \texttt{safetensors})
  \item \textbf{Timing detaliat} --- latență FP32 / FP16 / INT8 per-tensor /
    INT8 per-channel
\end{itemize}

% ===========================================================================
\begin{thebibliography}{9}
% ===========================================================================

\bibitem{dosovitskiy2021vit}
A. Dosovitskiy, L. Beyer, A. Kolesnikov et al.,
``An Image is Worth 16x16 Words: Transformers for Image Recognition at Scale,''
\emph{ICLR}, 2021. \href{https://arxiv.org/abs/2010.11929}{arXiv:2010.11929}.

\bibitem{touvron2021deit}
H. Touvron, M. Cord, M. Douze et al.,
``Training data-efficient image transformers \& distillation through attention,''
\emph{ICML}, 2021. \href{https://arxiv.org/abs/2012.12877}{arXiv:2012.12877}.

\bibitem{micikevicius2022fp8}
P. Micikevicius, D. Stosic, N. Burgess et al.,
``FP8 Formats for Deep Learning,''
2022. \href{https://arxiv.org/abs/2209.05433}{arXiv:2209.05433}.

\bibitem{jacob2018quantization}
B. Jacob, S. Kligys, B. Chen et al.,
``Quantization and Training of Neural Networks for Efficient Integer-Arithmetic-Only Inference,''
\emph{CVPR}, 2018. \href{https://arxiv.org/abs/1712.05877}{arXiv:1712.05877}.

\bibitem{imagenette}
J. Howard, ``Imagenette,'' fast.ai, 2019.
\url{https://github.com/fastai/imagenette}.

\bibitem{timm}
R. Wightman, ``PyTorch Image Models (timm),'' 2019.
\url{https://github.com/rwightman/pytorch-image-models}.

\bibitem{pytorch2023amp}
PyTorch Team,
``Automatic Mixed Precision package --- torch.amp,''
Documentație oficială PyTorch, 2023.
\url{https://pytorch.org/docs/stable/amp.html}.

\end{thebibliography}

\end{document}
